\setcounter{chapter}{11}
\pagestyle{fancy}

\chapter{重积分}


\section{二重积分}


\subsection{化为累次积分}

设 $f(x,y)$ 在 $xOy$ 平面有界闭区域 $D$ 上有定义, 且下面出现的积分都存在, 则:

\begin{enumerate}
    \item 当
          \[
              D_1=\{(x,y)\mid a\leq x\leq b,\ y_1(x)\leq y\leq y_2(x)\}
          \]
          (称为 $x$ 型区域)时,
          \[
              \iint\limits_{D_1}f(x,y)\,\mathrm{d}x\,\mathrm{d}y
              =\int_a^b\mathrm{d}x\int_{y_1(x)}^{y_2(x)}f(x,y)\,\mathrm{d}y.
          \]

    \item 当
          \[
              D_2=\{(x,y)\mid c\leq y\leq d,\ x_1(y)\leq x\leq x_2(y)\}
          \]
          (称为 $y$ 型区域)时,
          \[
              \iint\limits_{D_2}f(x,y)\,\mathrm{d}x\,\mathrm{d}y
              =\int_c^d\mathrm{d}y\int_{x_1(y)}^{x_2(y)}f(x,y)\,\mathrm{d}x.
          \]
\end{enumerate}

因此, 对于一般区域, 如区域可分成若干个不相交的 $x$ 型区域和 $y$ 型区域的并, 则可分别计算这些区域上积分的值再相加. 进一步, 也为后续的换元提供了方向----化为 $x$ 型区域(或 $y$ 型区域). 特别地, 若 $D=D_1=D_2$, 则上述两式均成立, 即两累次积分相等, 但具体计算时用哪个, 还得由被积函数决定.

\begin{example}
    设 $f(x,y)$ 在 $\mathbb{R}^2$ 上连续, $a>0$, 试交换下列累次积分的次序:
    \begin{enumerate}
        \item $\displaystyle\int_0^a\mathrm{d}y\int_{\frac{y}{2}}^{2y}f(x,y)\,\mathrm{d}x$;
        \item $\displaystyle\int_{-1}^1\mathrm{d}x\int_{x^2+x}^{x+1}f(x,y)\,\mathrm{d}y$;
        \item $\displaystyle\int_0^a\mathrm{d}y\int_{-y}^{\sqrt{\sqrt{y}}}f(x,y)\,\mathrm{d}x$;
        \item $\displaystyle\int_0^{2a}\mathrm{d}x\int_{\sqrt{2ax-x^2}}^{4ax}f(x,y)\,\mathrm{d}y$.
    \end{enumerate}
\end{example}
\exampleblank

\begin{solution}
    记积分区域为 $D$,直接按固定 $x$(或固定 $r$)作截线可得
    \begin{enumerate}
        \item
              \[
                  \int_0^{\frac{a}{2}}\!\mathrm{d}x\int_{\frac{x}{2}}^{2x}f(x,y)\,\mathrm{d}y
                  +\int_{\frac{a}{2}}^{2a}\!\mathrm{d}x\int_{\frac{x}{2}}^{a}f(x,y)\,\mathrm{d}y.
              \]
        \item 令 $x_{\pm}(y)=\dfrac{-1\pm\sqrt{1+4y}}{2}$,则原积分等于
              \[
                  \int_{-\frac{1}{4}}^{0}\!\mathrm{d}y\int_{x_-(y)}^{x_+(y)}f(x,y)\,\mathrm{d}x
                  +\int_0^2\!\mathrm{d}y\int_{y-1}^{x_+(y)}f(x,y)\,\mathrm{d}x.
              \]
        \item
              \[
                  \int_{-a}^{0}\!\mathrm{d}x\int_{-x}^{a}f(x,y)\,\mathrm{d}y
                  +\int_0^{\sqrt[4]{a}}\!\mathrm{d}x\int_{x^4}^{a}f(x,y)\,\mathrm{d}y.
              \]
        \item 题面给出的上下限在靠近 $x=0$ 时次序相反.按有向积分拆成两个正常区域,原式为
              \[
                  \int_0^{8a^2}\!\mathrm{d}y\int_{\frac{y}{4a}}^{2a}f(x,y)\,\mathrm{d}x
                  -\int_0^a\!\mathrm{d}y\int_{a-\sqrt{a^2-y^2}}^{a+\sqrt{a^2-y^2}}f(x,y)\,\mathrm{d}x.
              \]
    \end{enumerate}
\end{solution}

\begin{example}[\markstar]
    设 $f(x,y)$ 在 $\mathbb{R}^2$ 上连续, $a>0$, 试交换下列累次积分的次序:
    \begin{enumerate}
        \item $\displaystyle\int_0^{\frac{\pi}{2}}\mathrm{d}\theta\int_0^{2a\cos\theta}f(r,\theta)\,\mathrm{d}r$;
        \item $\displaystyle\int_0^{\frac{\pi}{4}}\mathrm{d}\theta\int_0^{a\sec\theta}f(r,\theta)\,\mathrm{d}r$;
        \item $\displaystyle\int_{-\frac{\pi}{2}}^{\frac{\pi}{2}}\mathrm{d}\theta\int_0^{2a\cos\theta}f(r\cos\theta,r\sin\theta)r\,\mathrm{d}r$.
    \end{enumerate}
\end{example}
\exampleblank

\begin{solution}
    按固定 $r$ 作截线,有
    \begin{enumerate}
        \item
              \[
                  \int_0^{2a}\!\mathrm{d}r\int_0^{\arccos\frac{r}{2a}}f(r,\theta)\,\mathrm{d}\theta.
              \]
        \item
              \[
                  \int_0^a\!\mathrm{d}r\int_0^{\frac{\pi}{4}}f(r,\theta)\,\mathrm{d}\theta
                  +\int_a^{\sqrt{2}a}\!\mathrm{d}r\int_{\arccos\frac{a}{r}}^{\frac{\pi}{4}}f(r,\theta)\,\mathrm{d}\theta.
              \]
        \item
              \[
                  \int_0^{2a}\!\mathrm{d}r
                  \int_{-\arccos\frac{r}{2a}}^{\arccos\frac{r}{2a}}
                  f(r\cos\theta,r\sin\theta)r\,\mathrm{d}\theta.
              \]
    \end{enumerate}
\end{solution}

\begin{example}
    证明下列等式:
    \begin{enumerate}
        \item
              \[
                  \int_0^e\mathrm{d}y\int_1^{y^2}\dfrac{\ln x}{e^x}\,\mathrm{d}x
                  +\int_e^{e^2}\mathrm{d}y\int_{\ln y}^{2}\dfrac{\ln x}{e^x}\,\mathrm{d}x
                  =\int_1^2\ln x\,\mathrm{d}x;
              \]
        \item
              \[
                  \int_a^b\mathrm{d}y\int_a^y(y-x)^n f(x)\,\mathrm{d}x
                  =\dfrac{1}{n+1}\int_a^b(b-x)^{n+1}f(x)\,\mathrm{d}x.
              \]
    \end{enumerate}
\end{example}
\exampleblank

\begin{solution}
    对于 (1),按题面直接交换积分,左端等于
    \[
        -\int_0^1\sqrt{x}\,\dfrac{\ln x}{e^x}\,\mathrm{d}x
        +\int_1^2\left(e^x-\sqrt{x}\right)\dfrac{\ln x}{e^x}\,\mathrm{d}x
        +\int_2^{e^2}\left(e-\sqrt{x}\right)\dfrac{\ln x}{e^x}\,\mathrm{d}x.
    \]
    其数值约为 $0.7421142902$,而右端为 $2\ln2-1\approx0.3862943611$,故 (1) 按现有题面不成立,原始资料中的积分限或被积函数需要核对.

    对于 (2),积分区域为 $a\leq x\leq y\leq b$.交换次序后
    \[
        \int_a^b f(x)\,\mathrm{d}x\int_x^b(y-x)^n\,\mathrm{d}y
        =\dfrac{1}{n+1}\int_a^b(b-x)^{n+1}f(x)\,\mathrm{d}x.
    \]
\end{solution}

\begin{theorem}[\markstar\ 对称性]
    \begin{enumerate}
        \item 积分区域 $D$ 关于 $x$ 轴对称.
              \begin{enumerate}[label=(\roman*)]
                  \item 若 $f(x,y)=-f(x,-y)$, 则 $\displaystyle\iint\limits_Df(x,y)\,\mathrm{d}x\,\mathrm{d}y=0$;
                  \item 若 $f(x,y)=f(x,-y)$, 则
                        \[
                            \iint\limits_Df(x,y)\,\mathrm{d}x\,\mathrm{d}y
                            =2\iint\limits_{D\cap\{y\geq0\}}f(x,y)\,\mathrm{d}x\,\mathrm{d}y.
                        \]
              \end{enumerate}

        \item 积分区域 $D$ 关于 $y$ 轴对称.
              \begin{enumerate}[label=(\roman*)]
                  \item 若 $f(x,y)=-f(-x,y)$, 则 $\displaystyle\iint\limits_Df(x,y)\,\mathrm{d}x\,\mathrm{d}y=0$;
                  \item 若 $f(x,y)=f(-x,y)$, 则
                        \[
                            \iint\limits_Df(x,y)\,\mathrm{d}x\,\mathrm{d}y
                            =2\iint\limits_{D\cap\{x\geq0\}}f(x,y)\,\mathrm{d}x\,\mathrm{d}y.
                        \]
              \end{enumerate}

        \item 积分区域 $D$ 关于原点对称.
              \begin{enumerate}[label=(\roman*)]
                  \item 若 $f(x,y)=-f(-x,-y)$, 则 $\displaystyle\iint\limits_Df(x,y)\,\mathrm{d}x\,\mathrm{d}y=0$;
                  \item 若 $f(x,y)=f(-x,-y)$, 则
                        \[
                            \iint\limits_Df(x,y)\,\mathrm{d}x\,\mathrm{d}y
                            =2\iint\limits_{D_1}f(x,y)\,\mathrm{d}x\,\mathrm{d}y,
                        \]
                        其中 $D_1$ 为区域 $D$ 的一半.
              \end{enumerate}
    \end{enumerate}
\end{theorem}

\begin{remark}
    对称性一般是用来简化计算, 消去部分被积函数项, 缩小积分区域.
\end{remark}

\begin{example}
    计算积分
    \[
        I=\iint\limits_Dx\left(1+yf(x^2+y^2)\right)\,\mathrm{d}x\,\mathrm{d}y,
    \]
    其中 $D$ 是由 $y=x^3,y=1,x=-1$ 所围成的区域, $f(x)$ 是实值连续函数.
\end{example}
\exampleblank

\begin{solution}
    写成累次积分并把关于 $x$ 的对称项配对:
    \[
        I=\int_{-1}^1\!\mathrm{d}x\int_{x^3}^1x\left(1+yf(x^2+y^2)\right)\,\mathrm{d}y.
    \]
    其中
    \[
        \int_{-1}^1x(1-x^3)\,\mathrm{d}x=-\dfrac{2}{5}.
    \]
    对任意 $x>0$,$x$ 与 $-x$ 对应的含 $f$ 部分之和为
    \[
        -x\int_{-x^3}^{x^3}y f(x^2+y^2)\,\mathrm{d}y=0.
    \]
    故 $I=-\dfrac{2}{5}$.
\end{solution}

\begin{example}
    计算积分:
    \begin{enumerate}
        \item
              \[
                  \iint\limits_D\left|xy-\dfrac{1}{4}\right|\,\mathrm{d}x\,\mathrm{d}y,
                  \qquad D=\lbrack0,1\rbrack\times\lbrack0,1\rbrack;
              \]
        \item
              \[
                  \iint\limits_{|x|\leq1,\ 0\leq y\leq2}\sqrt{|y-x^2|}\,\mathrm{d}x\,\mathrm{d}y.
              \]
    \end{enumerate}
\end{example}
\exampleblank

\begin{solution}
    分别按绝对值变号曲线分区.
    \begin{enumerate}
        \item 对 $x\leq\dfrac{1}{4}$ 不变号,对 $x>\dfrac{1}{4}$ 在 $y=\dfrac{1}{4x}$ 处分开,得到
              \[
                  \iint_D\left|xy-\dfrac{1}{4}\right|\,\mathrm{d}x\,\mathrm{d}y
                  =\dfrac{3}{32}+\dfrac{\ln2}{8}.
              \]
        \item 利用关于 $y$ 轴的对称性,先对 $y$ 积分:
              \[
                  2\int_0^1\left[\int_0^{x^2}\sqrt{x^2-y}\,\mathrm{d}y
                      +\int_{x^2}^2\sqrt{y-x^2}\,\mathrm{d}y\right]\mathrm{d}x
                  =\dfrac{5}{3}+\dfrac{\pi}{2}.
              \]
    \end{enumerate}
\end{solution}

\begin{example}
    计算积分:
    \begin{enumerate}
        \item $\displaystyle\iint\limits_{x^2+y^2\leq5}\operatorname{sgn}(x^2-y^2+3)\,\mathrm{d}x\,\mathrm{d}y$;
        \item $\displaystyle\iint\limits_{0\leq x,y\leq2}[x+y]\,\mathrm{d}x\,\mathrm{d}y$.
    \end{enumerate}
\end{example}
\exampleblank

\begin{solution}
    \begin{enumerate}
        \item 负号区域由 $y^2-x^2>3$ 与 $x^2+y^2\leq5$ 给出.计算其面积并用``正面积减负面积'',得
              \[
                  -5\pi+20\arcsin\dfrac{2}{\sqrt5}+6\ln3.
              \]
        \item 令 $s=x+y$.正方形中截线长度在 $0\leq s\leq2$ 时为 $s$,在 $2\leq s\leq4$ 时为 $4-s$,故
              \[
                  \int_0^2[s]s\,\mathrm{d}s+\int_2^4[s](4-s)\,\mathrm{d}s=6.
              \]
    \end{enumerate}
\end{solution}

\begin{example}
    举例说明下列命题:
    \begin{enumerate}
        \item 设 $D=\{(x,y)\mid a\leq x\leq b,\ c<y\leq d\}$, $f\in C(D)$, 此时累次积分不一定能交换次序;
        \item 存在 $\lbrack0,1\rbrack\times\lbrack0,1\rbrack$ 上非负不可积函数 $f(x,y)$, 使得
              \[
                  \int_0^1\mathrm{d}x\int_0^1f(x,y)\,\mathrm{d}y
                  =0
                  =\int_0^1\mathrm{d}y\int_0^1f(x,y)\,\mathrm{d}x.
              \]
    \end{enumerate}
\end{example}
\exampleblank

\begin{solution}
    \begin{enumerate}
        \item 取 $D=\{(x,y)\mid0\leq x\leq1,\ 0<y\leq1\}$ 及
              \[
                  f(x,y)=\dfrac{xy(x^2-y^2)}{(x^2+y^2)^3}.
              \]
              它在 $D$ 上连续,并且
              \[
                  \int_0^1\!\mathrm{d}x\int_0^1f(x,y)\,\mathrm{d}y=\dfrac{1}{8},
                  \qquad
                  \int_0^1\!\mathrm{d}y\int_0^1f(x,y)\,\mathrm{d}x=-\dfrac{1}{8}.
              \]
        \item 取一个图像在单位正方形中稠密的双射 $g:[0,1]\to[0,1]$,令 $f$ 为其图像的示性函数.每条水平线和竖直线与图像都只有一个交点,故两个累次 Riemann 积分均为零;但图像及其补集都稠密,$f$ 处处不连续,因而在正方形上不可积.
    \end{enumerate}
\end{solution}



\subsection{变量替换}

\begin{theorem}[\markstar]
    设 $D,\widetilde D$ 是 $\mathbb{R}^2$ 中由逐段光滑曲线围成的有界闭区域, 变换
    \[
        T:\begin{cases}
            x=x(u,v), \\
            y=y(u,v)
        \end{cases}
        :\widetilde D\to D
    \]
    是同胚变换(双射及 $T,T^{-1}$ 连续), 且 $T\in C^{(2)}(\widetilde D)$, 记
    \[
        J=\dfrac{\partial(x,y)}{\partial(u,v)}.
    \]
    若 $f(x,y)$ 在 $D$ 上可积, 则
    \[
        \iint\limits_Df(x,y)\,\mathrm{d}x\,\mathrm{d}y
        =\iint\limits_{\widetilde D}f\left(x(u,v),y(u,v)\right)|J(u,v)|\,\mathrm{d}u\,\mathrm{d}v.
    \]
\end{theorem}

\begin{remark}
    特别地,
    \[
        \begin{cases}
            x=r\cos\theta, \\
            y=r\sin\theta
        \end{cases}
    \]
    为极坐标变换, $\theta$ 为沿 $x$ 正半轴逆时针方向的旋转角, $r$ 为到原点的距离, 此时 $J=r$. 换元的目的旨在化简积分区域(化为 $x$ 型或 $y$ 型或二者多个的并)和被积函数.
\end{remark}

\begin{example}
    计算积分:
    \begin{enumerate}
        \item
              \[
                  I=\iint\limits_D\dfrac{3x}{y^2+xy}\,\mathrm{d}x\,\mathrm{d}y,
              \]
              其中 $D$ 为平面曲线 $xy=1,xy=3,y^2=x,y^2=3x$ 所围成的区域;
        \item
              \[
                  J=\iint\limits_D\dfrac{(\sqrt{x}+\sqrt{y})^4}{x^2}\,\mathrm{d}x\,\mathrm{d}y,
              \]
              其中 $D$ 为由 $x$ 轴,$y=x$,$\sqrt{x}+\sqrt{y}=1$ 和 $\sqrt{x}+\sqrt{y}=2$ 所围成的区域.
    \end{enumerate}
\end{example}
\exampleblank

\begin{solution}
    \begin{enumerate}
        \item 令 $u=xy=p^3$,$v=y^2/x=q^3$.在所围的第一象限区域中 $1\leq p,q\leq\sqrt[3]{3}$,且
              \[
                  x=\dfrac{p^2}{q},\qquad y=pq,\qquad
                  \left|\dfrac{\partial(x,y)}{\partial(p,q)}\right|=\dfrac{3p^2}{q}.
              \]
              记 $c=\sqrt[3]{3}$ 及
              \[
                  \Phi_a(q)=\dfrac{q^3}{3}\ln(q^2+a)+\dfrac{2aq}{3}
                  -\dfrac{2a^{3/2}}{3}\arctan\dfrac{q}{\sqrt a},
              \]
              则积分的精确值为
              \[
                  9\left[-\dfrac{c^2-1}{2q}-(c-1)q+\Phi_c(q)-\Phi_1(q)\right]_{q=1}^{q=c}.
              \]
        \item 令 $u=\sqrt{x}+\sqrt{y}$,$t=\sqrt{y/x}$.则 $1\leq u\leq2$,$0\leq t\leq1$,且被积式连同 Jacobian 化为 $4u^3t$,所以
              \[
                  J=\int_1^2\!\int_0^1 4u^3t\,\mathrm{d}t\,\mathrm{d}u=\dfrac{15}{2}.
              \]
    \end{enumerate}
\end{solution}

\begin{example}
    计算积分:
    \begin{enumerate}
        \item
              \[
                  \iint\limits_D\dfrac{x^2-y^2}{\sqrt{x+y+3}}\,\mathrm{d}x\,\mathrm{d}y,
                  \qquad D=\{(x,y)\mid |x|+|y|\leq1\};
              \]
        \item
              \[
                  \iint\limits_D\dfrac{(x+y)\ln(1+xy)}{\sqrt{1-x-y}}\,\mathrm{d}x\,\mathrm{d}y,
                  \qquad D=\left\{(x,y)\mid0\leq y\leq x,\ \dfrac{3}{4}\leq x+y\leq1\right\};
              \]
        \item
              \[
                  \iint\limits_De^{\frac{x-y}{x+y}}\,\mathrm{d}x\,\mathrm{d}y,
              \]
              其中 $D$ 是 $x=0,y=0,x+y=1$ 所围成的有界闭区域.
    \end{enumerate}
\end{example}
\exampleblank

\begin{solution}
    \begin{enumerate}
        \item 区域关于 $x,y$ 互换不变,而被积函数在该变换下变号,故积分为 $0$.
        \item 令 $u=x+y$,$v=x-y$,则 $\dfrac{3}{4}\leq u\leq1$,$0\leq v\leq u$,$\mathrm{d}x\,\mathrm{d}y=\dfrac{1}{2}\mathrm{d}u\,\mathrm{d}v$.先积去 $v$ 后得到精确的一维表示
              \[
                  \int_{\frac{3}{4}}^{1}\dfrac{u}{\sqrt{1-u}}
                  \left[\sqrt{4+u^2}\operatorname{arctanh}\dfrac{u}{\sqrt{4+u^2}}-u\right]\mathrm{d}u.
              \]
        \item 令 $u=x+y$,$v=(x-y)/(x+y)$,则 $0\leq u\leq1$,$-1\leq v\leq1$,Jacobian 为 $u/2$.故
              \[
                  \iint_D e^{\frac{x-y}{x+y}}\,\mathrm{d}x\,\mathrm{d}y
                  =\dfrac{1}{4}\int_{-1}^1e^v\,\mathrm{d}v
                  =\dfrac{1}{4}\left(e-e^{-1}\right).
              \]
    \end{enumerate}
\end{solution}

\begin{example}
    计算积分:
    \begin{enumerate}
        \item $\displaystyle\iint\limits_D(x+y)\,\mathrm{d}x\,\mathrm{d}y$, $D$ 是由曲线 $x^2+y^2=x+y$ 所围成的区域;
        \item
              \[
                  \iint\limits_D\dfrac{1}{xy}\,\mathrm{d}x\,\mathrm{d}y,
              \]
              其中 $D$ 为平面上由
              \[
                  2\leq\dfrac{x}{x^2+y^2},\qquad
                  \dfrac{y}{x^2+y^2}\leq4
              \]
              确定的区域;
        \item
              \[
                  \iint\limits_D\left(3x^3+x^2+y^2+2x-2y+1\right)\,\mathrm{d}x\,\mathrm{d}y,
              \]
              其中
              \[
                  D=\{(x,y)\mid1\leq x^2+(y-1)^2\leq2,\ x^2+y^2\leq1\}.
              \]
    \end{enumerate}
\end{example}
\exampleblank

\begin{solution}
    \begin{enumerate}
        \item 圆域的圆心为 $(1/2,1/2)$,面积为 $\pi/2$.由形心公式,积分为 $\pi/2$.
        \item 按现有题面,反演
              \[
                  u=\dfrac{x}{x^2+y^2},\qquad v=\dfrac{y}{x^2+y^2}
              \]
              把区域变为 $u\geq2,v\leq4$,且被积式连同 Jacobian 为 $\mathrm{d}u\,\mathrm{d}v/(uv)$.该区域穿过 $v=0$ 且无下界,积分发散.因此题面若预期有限值,尚缺少关于 $v$(或原变量)的边界条件.
        \item 以 $(0,1)$ 为极点作极坐标,关于 $x$ 的奇函数项积分为零,故
              \[
                  I=\int_1^{\sqrt2}\rho^3\left(\pi-2\arcsin\dfrac{\rho}{2}\right)\mathrm{d}\rho
                  =-2+\dfrac{7\sqrt3}{8}+\dfrac{7\pi}{12}.
              \]
    \end{enumerate}
\end{solution}

\begin{example}
    计算积分:
    \begin{enumerate}
        \item
              \[
                  \iint\limits_{x^2+y^2\leq\frac{3}{16}}
                  \min\left\{\sqrt[3]{\dfrac{3}{16}-x^2-y^2},\ 2(x^2+y^2)\right\}\,\mathrm{d}x\,\mathrm{d}y;
              \]
        \item
              \[
                  \iint\limits_{x^2+y^2\leq1}\left|\dfrac{x+y}{\sqrt{2}}-x^2-y^2\right|\,\mathrm{d}x\,\mathrm{d}y.
              \]
    \end{enumerate}
\end{example}
\exampleblank

\begin{solution}
    \begin{enumerate}
        \item 令 $s=r^2$,并令 $\alpha\in(0,3/16)$ 为方程
              \[
                  8\alpha^3+\alpha=\dfrac{3}{16}
              \]
              的唯一根.两项在 $s=\alpha$ 时相等,故原积分为
              \[
                  \pi\left[\int_0^\alpha2s\,\mathrm{d}s
                      +\int_\alpha^{\frac{3}{16}}\left(\dfrac{3}{16}-s\right)^{\frac{1}{3}}\mathrm{d}s\right]
                  =\pi\left(\alpha^2+12\alpha^4\right).
              \]
        \item 旋转坐标后,被积函数为 $|r\cos\theta-r^2|$.对角变量先积分,有
              \[
                  \int_0^{2\pi}|\cos\theta-r|\,\mathrm{d}\theta
                  =4\sqrt{1-r^2}+2\pi r-4r\arccos r.
              \]
              再对 $0\leq r\leq1$ 积分,结果为 $9\pi/16$.
    \end{enumerate}
\end{solution}

\begin{example}
    计算积分:
    \begin{enumerate}
        \item $\displaystyle\iint\limits_{x^2+4y^2\leq1}(x^2+y^2)\,\mathrm{d}x\,\mathrm{d}y$;
        \item $\displaystyle\iint\limits_{\sqrt{x}+\sqrt{y}\leq1,\ x,y\geq0}xy\,\mathrm{d}x\,\mathrm{d}y$;
        \item $\displaystyle\iint\limits_{x^4+y^4\leq1}(x^2+y^2)\,\mathrm{d}x\,\mathrm{d}y$.
    \end{enumerate}
\end{example}
\exampleblank

\begin{solution}
    \begin{enumerate}
        \item 作 $x=u,y=v/2$,化到单位圆后由对称性得到 $5\pi/32$.
        \item 令 $u=\sqrt{x},v=\sqrt{y}$.则
              \[
                  4\int_{u+v\leq1}u^3v^3\,\mathrm{d}u\,\mathrm{d}v=\dfrac{1}{280}.
              \]
        \item 由四象限对称性和 Dirichlet 积分,
              \[
                  \iint_{x^4+y^4\leq1}(x^2+y^2)\,\mathrm{d}x\,\mathrm{d}y
                  =\dfrac{1}{2}\Gamma\left(\dfrac{1}{4}\right)\Gamma\left(\dfrac{3}{4}\right)
                  =\dfrac{\pi}{\sqrt2}.
              \]
    \end{enumerate}
\end{solution}

\begin{example}
    计算积分
    \[
        \iint\limits_D\sqrt{\dfrac{xy}{x+y}}\,\mathrm{d}x\,\mathrm{d}y,
    \]
    其中 $D$ 是由 $x=0,y=0,x+y=1$ 所围成的区域.
\end{example}
\exampleblank

\begin{solution}
    令 $u=x+y$,$v=x/(x+y)$,则 $0\leq u,v\leq1$,$x=uv$,$y=u(1-v)$,Jacobian 为 $u$.于是
    \[
        I=\int_0^1u^{\frac{3}{2}}\,\mathrm{d}u
        \int_0^1v^{\frac{1}{2}}(1-v)^{\frac{1}{2}}\,\mathrm{d}v
        =\dfrac{2}{5} B\left(\dfrac{3}{2},\dfrac{3}{2}\right)=\dfrac{\pi}{20}.
    \]
\end{solution}



\subsection{理论证明}

由一元函数积分中值定理证明过程易知, 重积分中值定理仍成立.

\begin{theorem}[\markstar]
    设 $f(x,y)$ 在有界闭区域 $D$ 上连续, 则存在 $(\xi,\eta)\in D$, 使得
    \[
        \iint\limits_Df(x,y)\,\mathrm{d}x\,\mathrm{d}y=f(\xi,\eta)|D|.
    \]
    其中 $|D|$ 表示 $D$ 的面积.
\end{theorem}

\begin{remark}
    事实上可以证明上述中值点 $(\xi,\eta)$ 可以有无穷多个.
\end{remark}

\begin{example}
    设 $f(x,y)\in C^{(2)}(D)$, $D=\lbrack a,b\rbrack\times\lbrack c,d\rbrack$, 则
    \begin{enumerate}
        \item
              \[
                  \iint\limits_D f_{xy}(x,y)\,\mathrm{d}x\,\mathrm{d}y
                  =\iint\limits_D f_{yx}(x,y)\,\mathrm{d}x\,\mathrm{d}y;
              \]
        \item 利用 (1) 证明 $f_{xy}(x,y)=f_{yx}(x,y)$.
    \end{enumerate}
\end{example}
\exampleblank

\begin{solution}
    由微积分基本定理,
    \[
        \iint_Df_{xy}\,\mathrm{d}x\,\mathrm{d}y
        =f(b,d)-f(a,d)-f(b,c)+f(a,c),
    \]
    而先按相反次序积分,$f_{yx}$ 的积分也是同一边界差,故 (1) 成立.

    把上述结论用于任意子矩形 $[x,x+h]\times[y,y+k]\subset D$,可得
    \[
        \iint_{[x,x+h]\times[y,y+k]}(f_{xy}-f_{yx})\,\mathrm{d}x\,\mathrm{d}y=0.
    \]
    除以 $hk$ 并令 $h,k\to0$,利用连续性即得 $f_{xy}(x,y)=f_{yx}(x,y)$.
\end{solution}

\begin{example}
    设 $f\in C(\mathbb{R})$, $h>0$,
    \[
        F(x)=\dfrac{1}{h^2}\int_0^h\mathrm{d}y\int_0^h f(x+y+z)\,\mathrm{d}z,
    \]
    求 $F''(x)$.
\end{example}
\exampleblank

\begin{solution}
    先把内层积分写成差分:
    \[
        F'(x)=\dfrac{1}{h^2}\int_0^h\left[f(x+y+h)-f(x+y)\right]\,\mathrm{d}y.
    \]
    再次使用微积分基本定理,得到
    \[
        F''(x)=\dfrac{f(x+2h)-2f(x+h)+f(x)}{h^2}.
    \]
\end{solution}

\begin{example}
    设 $x=x(u,v),y=y(u,v)$ 有连续偏导数, 一一对应地将区域 $D'$ 映射到 $xOy$ 平面的区域 $D$, 满足 $J\ne0$, 且
    \[
        \dfrac{\partial x}{\partial u}=\dfrac{\partial y}{\partial v},\qquad
        \dfrac{\partial x}{\partial v}=-\dfrac{\partial y}{\partial u},
    \]
    试证
    \[
        \iint\limits_D\left[\left(\dfrac{\partial f}{\partial x}\right)^2+
            \left(\dfrac{\partial f}{\partial y}\right)^2\right]\mathrm{d}x\,\mathrm{d}y
        =\iint\limits_{D'}\left[\left(\dfrac{\partial f}{\partial u}\right)^2+
            \left(\dfrac{\partial f}{\partial v}\right)^2\right]\mathrm{d}u\,\mathrm{d}v.
    \]
\end{example}
\exampleblank

\begin{solution}
    记 $a=x_u=y_v$,$b=x_v=-y_u$,则 Jacobian 的绝对值为 $a^2+b^2$.链式法则给出
    \[
        f_u^2+f_v^2=(a^2+b^2)(f_x^2+f_y^2).
    \]
    另一方面,$\mathrm{d}x\,\mathrm{d}y=(a^2+b^2)\mathrm{d}u\,\mathrm{d}v$.代入变量替换公式即得所证等式.
\end{solution}

\begin{example}
    设 $f(x,y)\geq0$ 在 $D:x^2+y^2\leq a^2$ 上有连续的一阶偏导数, 边界上取值为零, 证明
    \[
        \left|\iint\limits_Df(x,y)\,\mathrm{d}x\,\mathrm{d}y\right|
        \leq\dfrac{\pi a^3}{3}
        \max_{(x,y)\in D}\sqrt{\left(\dfrac{\partial f}{\partial x}\right)^2+
            \left(\dfrac{\partial f}{\partial y}\right)^2}.
    \]
\end{example}
\exampleblank

\begin{solution}
    记右端最大值为 $M$.由于 $f$ 在边界上为零,对极坐标点 $(r,\theta)$ 沿半径使用中值定理,
    \[
        0\leq f(r,\theta)\leq M(a-r).
    \]
    所以
    \[
        \left|\iint_Df\right|\leq2\pi M\int_0^a(a-r)r\,\mathrm{d}r
        =\dfrac{\pi a^3}{3}M.
    \]
\end{solution}

\begin{example}
    设 $D=\{(x,y)\mid0\leq x,y\leq1\}$, $f(x,y)$ 在 $D$ 上四次连续可微, 当 $(x,y)\in\partial D$ 时 $f(x,y)=0$. 对任意 $(x,y)\in D$, 有
    \[
        \left|\dfrac{\partial^4f(x,y)}{\partial x^2\partial y^2}\right|\leq M,
    \]
    则
    \[
        \left|\iint\limits_Df(x,y)\,\mathrm{d}x\,\mathrm{d}y\right|\leq\dfrac{M}{144}.
    \]
\end{example}
\exampleblank

\begin{solution}
    若 $g(0)=g(1)=0$,两次分部积分给出
    \[
        \int_0^1g(x)\,\mathrm{d}x=-\dfrac{1}{2}\int_0^1x(1-x)g''(x)\,\mathrm{d}x.
    \]
    分别对 $x,y$ 使用此式,得到
    \[
        \iint_Df=\dfrac{1}{4}\iint_Dx(1-x)y(1-y)f_{xxyy}(x,y)\,\mathrm{d}x\,\mathrm{d}y.
    \]
    因此
    \[
        \left|\iint_Df\right|\leq\dfrac{M}{4}\left(\int_0^1x(1-x)\,\mathrm{d}x\right)^2
        =\dfrac{M}{144}.
    \]
\end{solution}

\begin{example}[\markstar]
    利用高维积分证明低维积分不等式:
    \begin{enumerate}
        \item (Minkowski 不等式)设 $f(x,y)$ 在 $\mathbb{R}^2$ 上连续, 则
              \[
                  \left\{\int_a^b\left[\int_c^d f(x,y)\,\mathrm{d}y\right]^2\mathrm{d}x\right\}^{\frac{1}{2}}
                  \leq
                  \int_c^d\left\{\int_a^b f^2(x,y)\,\mathrm{d}x\right\}^{\frac{1}{2}}\mathrm{d}y.
              \]
        \item 设 $f(x)$ 是 $\lbrack0,1\rbrack$ 上递减正值函数, 则
              \[
                  \dfrac{\displaystyle\int_0^1x f^2(x)\,\mathrm{d}x}
                  {\displaystyle\int_0^1x f(x)\,\mathrm{d}x}
                  \leq
                  \dfrac{\displaystyle\int_0^1 f^2(x)\,\mathrm{d}x}
                  {\displaystyle\int_0^1 f(x)\,\mathrm{d}x}.
              \]
        \item 设 $p(x)\in R\lbrack a,b\rbrack$, $p(x)\geq0$, $f(x),g(x)$ 在 $\lbrack a,b\rbrack$ 上递增, 则
              \[
                  \left(\int_a^bp(x)f(x)\,\mathrm{d}x\right)
                  \left(\int_a^bp(x)g(x)\,\mathrm{d}x\right)
                  \leq
                  \left(\int_a^bp(x)\,\mathrm{d}x\right)
                  \left(\int_a^bp(x)f(x)g(x)\,\mathrm{d}x\right).
              \]
    \end{enumerate}
\end{example}
\exampleblank

\begin{solution}
    \begin{enumerate}
        \item 展开左端平方并使用 Cauchy--Schwarz 不等式:
              \[
                  \begin{aligned}
                      \int_a^b\left(\int_c^df(x,y)\,\mathrm{d}y\right)^2\mathrm{d}x
                       & =\int_c^d\!\int_c^d\!\int_a^b f(x,y)f(x,z)\,\mathrm{d}x\,\mathrm{d}y\,\mathrm{d}z       \\
                       & \leq\left[\int_c^d\left(\int_a^bf(x,y)^2\,\mathrm{d}x\right)^{1/2}\mathrm{d}y\right]^2.
                  \end{aligned}
              \]
        \item 交叉相乘后的右端减左端等于
              \[
                  \dfrac{1}{2}\int_0^1\!\int_0^1 f(x)f(y)(y-x)(f(x)-f(y))\,\mathrm{d}x\,\mathrm{d}y\geq0,
              \]
              因为 $f$ 递减.
        \item 两边之差为
              \[
                  \dfrac{1}{2}\int_a^b\!\int_a^bp(x)p(y)(f(x)-f(y))(g(x)-g(y))\,\mathrm{d}x\,\mathrm{d}y\geq0.
              \]
    \end{enumerate}
\end{solution}



\subsection{二重积分的应用}

\begin{example}
    求由下列曲面围成的立体体积 $V$:
    \begin{enumerate}
        \item $z=0,z=x^2+y^2,x^2+y^2=x,x^2+y^2=2x$;
        \item $z=x^2+2y^2,z=2-x^2$;
        \item $\dfrac{x^2}{a^2}+\dfrac{y^2}{b^2}+\dfrac{z^2}{c^2}=1$, $\dfrac{x^2}{a^2}+\dfrac{y^2}{b^2}=\dfrac{z^2}{c^2}$, $z>0$, $a,b,c>0$.
    \end{enumerate}
\end{example}
\exampleblank

\begin{solution}
    \begin{enumerate}
        \item 极坐标中 $\cos\theta\leq r\leq2\cos\theta$,$-\pi/2\leq\theta\leq\pi/2$,故
              \[
                  V=\dfrac{1}{4}\int_{-\frac{\pi}{2}}^{\frac{\pi}{2}}\left(16-1\right)\cos^4\theta\,\mathrm{d}\theta
                  =\dfrac{45\pi}{32}.
              \]
        \item 投影为单位圆,竖直高度为 $2(1-x^2-y^2)$,故 $V=\pi$.
        \item 作 $x=aX,y=bY,z=cZ$.若所求为椭球内且圆锥上方的通常封闭区域,则
              \[
                  V=abc\int_0^{2\pi}\!\int_0^{\frac{\pi}{4}}\!\int_0^1r^2\sin\varphi\,\mathrm{d}r\,\mathrm{d}\varphi\,\mathrm{d}\theta
                  =\dfrac{2\pi abc}{3}\left(1-\dfrac{1}{\sqrt{2}}\right).
              \]
    \end{enumerate}
\end{solution}

\begin{example}
    求下列曲线围成的区域面积 $S$:
    \begin{enumerate}
        \item $x=0,y=0,\sqrt[4]{\dfrac{x}{a}}+\sqrt[4]{\dfrac{y}{b}}=1$, $a,b>0$;
        \item $\left(\dfrac{x^2}{a^2}+\dfrac{y^2}{b^2}\right)^2=\dfrac{xy}{c^2}$, $a,b,c>0$;
        \item $\dfrac{x^2}{a^2}+\dfrac{y^2}{b^2}=hx+ky$, $a,b,h,k>0$.
    \end{enumerate}
\end{example}
\exampleblank

\begin{solution}
    \begin{enumerate}
        \item 令 $u=\sqrt[4]{x/a}$,则 $y=b(1-u)^4$,从而
              \[
                  S=4ab\int_0^1u^3(1-u)^4\,\mathrm{d}u=\dfrac{ab}{70}.
              \]
        \item 极坐标中曲线只位于第一,第三象限,直接积分得
              \[
                  S=\dfrac{1}{c^2}\int_0^{\frac{\pi}{2}}
                  \dfrac{\sin\theta\cos\theta}{(\cos^2\theta/a^2+\sin^2\theta/b^2)^2}\,\mathrm{d}\theta
                  =\dfrac{a^2b^2}{2c^2}.
              \]
        \item 配方后所得椭圆的归一化半径平方为 $(a^2h^2+b^2k^2)/4$,故
              \[
                  S=\dfrac{\pi ab}{4}(a^2h^2+b^2k^2).
              \]
    \end{enumerate}
\end{solution}

\begin{example}[\markstar]
    用二重积分计算
    \[
        \int_0^{+\infty}e^{-x^2}\,\mathrm{d}x.
    \]
\end{example}
\exampleblank

\begin{solution}
    令 $I=\int_0^{+\infty}e^{-x^2}\,\mathrm{d}x$.由 Tonelli 定理与极坐标变换,
    \[
        I^2=\int_0^{\frac{\pi}{2}}\!\mathrm{d}\theta\int_0^{+\infty}e^{-r^2}r\,\mathrm{d}r
        =\dfrac{\pi}{4}.
    \]
    因 $I>0$,故 $I=\sqrt\pi/2$.
\end{solution}


\newpage

\section{三重积分}

本节整体内容与上节差不多, 分为重积分化累次积分,变量替换,理论证明以及应用, 但大部分并没有什么新的东西, 只不过升了一维导致计算更加复杂了些. 值得一提的是, 本节变量替换中多了正交变换与柱坐标变换.



\subsection{化为累次积分}

不同于二重积分的换序, 三重积分的换序涉及到立体图形, 可能会因图像复杂而难以刻画积分区域, 进而难以换序. 事实上, 可以换个思路, 不必 $x,y,z$ 三个变量一起换, 可以视 $z$ 为常数, 然后先只换 $x,y$, 此即大大降低了难度, 进而可通过多换几次二元序达到换三元序的目的.

\begin{example}[\markstar]
    试将下列三重积分次序 $z\to y\to x$ 改变为 $x\to y\to z$:
    \begin{enumerate}
        \item $\displaystyle\int_0^1\mathrm{d}x\int_0^{1-x}\mathrm{d}y\int_0^{x+y}f(x,y,z)\,\mathrm{d}z$;
        \item $\displaystyle\int_{-1}^1\mathrm{d}x\int_{-\sqrt{1-x^2}}^{\sqrt{1-x^2}}\mathrm{d}y\int_{\sqrt{x^2+y^2}}^1f(x,y,z)\,\mathrm{d}z$;
        \item $\displaystyle\int_0^1\mathrm{d}x\int_0^x\mathrm{d}y\int_0^{xy}f(x,y,z)\,\mathrm{d}z$.
    \end{enumerate}
\end{example}
\exampleblank

\begin{solution}
    按 $z$ 为最外层变量描述区域:
    \begin{enumerate}
        \item
              \[
                  \int_0^1\!\mathrm{d}z\left[
                      \int_0^z\!\mathrm{d}y\int_{z-y}^{1-y}f\,\mathrm{d}x
                      +\int_z^1\!\mathrm{d}y\int_0^{1-y}f\,\mathrm{d}x\right].
              \]
        \item
              \[
                  \int_0^1\!\mathrm{d}z\int_{-z}^{z}\!\mathrm{d}y
                  \int_{-\sqrt{z^2-y^2}}^{\sqrt{z^2-y^2}}f\,\mathrm{d}x.
              \]
        \item
              \[
                  \int_0^1\!\mathrm{d}z\left[
                      \int_z^{\sqrt z}\!\mathrm{d}y\int_{z/y}^{1}f\,\mathrm{d}x
                      +\int_{\sqrt z}^{1}\!\mathrm{d}y\int_y^1f\,\mathrm{d}x\right].
              \]
    \end{enumerate}
\end{solution}

\begin{theorem}[\markstar\ 对称性]
    \begin{enumerate}
        \item 积分区域 $V$ 关于 $xy$ 平面对称, 则考虑 $f(x,y,z)$ 关于 $z$ 的奇偶性;
        \item 积分区域 $V$ 关于 $yz$ 平面对称, 则考虑 $f(x,y,z)$ 关于 $x$ 的奇偶性;
        \item 积分区域 $V$ 关于 $xz$ 平面对称, 则考虑 $f(x,y,z)$ 关于 $y$ 的奇偶性;
        \item 积分区域 $V$ 关于原点对称, 则考虑 $f(x,y,z)$ 关于原点的对称性.
    \end{enumerate}
    相应地, 奇对称项在对称区域上的三重积分为 $0$, 偶对称项的积分可化为半区域积分的 $2$ 倍.
\end{theorem}

\begin{example}
    计算积分:
    \begin{enumerate}
        \item
              \[
                  \iiint\limits_V\dfrac{z\ln(x^2+y^2+z^2+1)}{x^2+y^2+z^2+1}\,\mathrm{d}x\,\mathrm{d}y\,\mathrm{d}z,
              \]
              其中 $V:x^2+y^2+z^2\leq1$;
        \item
              \[
                  \iiint\limits_{\substack{x^2+y^2+z^2\leq1\\x^2+y^2-z^2\geq\frac{1}{2}}}z\,\mathrm{d}x\,\mathrm{d}y\,\mathrm{d}z.
              \]
    \end{enumerate}
\end{example}
\exampleblank

\begin{solution}
    两个积分区域都关于平面 $z=0$ 对称,而被积函数关于 $z$ 为奇函数,故两积分均为 $0$.
\end{solution}

\begin{example}
    求积分
    \[
        \iiint\limits_V(x+y+z)\,\mathrm{d}x\,\mathrm{d}y\,\mathrm{d}z,
    \]
    其中 $V$ 是由平面 $x+y+z=1$ 以及三个坐标平面所围成的区域.
\end{example}
\exampleblank

\begin{solution}
    由对称性,三个坐标的一阶矩相同.四面体体积为 $1/6$,形心为 $(1/4,1/4,1/4)$,所以
    \[
        \iiint_V(x+y+z)\,\mathrm{d}V
        =\dfrac{1}{6}\left(\dfrac{1}{4}+\dfrac{1}{4}+\dfrac{1}{4}\right)=\dfrac{1}{8}.
    \]
\end{solution}

\begin{example}
    求区域
    \[
        V:\quad0\leq x\leq1,\quad0\leq y\leq x,\quad x+y\leq z\leq e^{x+y}
    \]
    的体积.
\end{example}
\exampleblank

\begin{solution}
    直接对 $z$ 积分:
    \[
        V=\int_0^1\!\mathrm{d}x\int_0^x\left(e^{x+y}-x-y\right)\,\mathrm{d}y
        =\dfrac{(e-1)^2}{2}-\dfrac{1}{2}
        =\dfrac{e(e-2)}{2}.
    \]
\end{solution}

\begin{example}
    设 $\Omega$ 由 $z^2=x^2+y^2,z=0,xy=1,xy=2,y=3x,y=4x$ 所围成, 求积分
    \[
        \iiint\limits_\Omega x^2y^2z\,\mathrm{d}x\,\mathrm{d}y\,\mathrm{d}z.
    \]
\end{example}
\exampleblank

\begin{solution}
    取 $u=xy$,$v=y/x$.在第一象限分支上 Jacobian 为 $1/(2v)$,第三象限给出相同贡献.先对 $z$ 积分后,
    \[
        I=\dfrac{1}{2}\int_1^2u^3\,\mathrm{d}u\int_3^4\left(1+\dfrac{1}{v^2}\right)\mathrm{d}v
        =\dfrac{65}{32}.
    \]
    这里按六个曲面围成的两个对称分支都属于 $\Omega$ 理解;若原题只取第一象限分支,答案应为 $65/64$.
\end{solution}

\begin{example}
    计算积分
    \[
        \iiint\limits_V\dfrac{\mathrm{d}V}{y^2+z^2},
    \]
    其中 $V$ 为一棱台, 其六个顶点为 $A(0,0,1)$, $B(0,1,1)$, $C(1,1,1)$, $D(0,0,2)$, $E(0,2,2)$, $F(2,2,2)$.
\end{example}
\exampleblank

\begin{solution}
    在高度 $z\in[1,2]$ 的截面上,三角形满足 $0\leq x\leq y\leq z$.故
    \[
        I=\int_1^2\!\mathrm{d}z\int_0^z\!\mathrm{d}y\int_0^y\dfrac{\mathrm{d}x}{y^2+z^2}
        =\int_1^2\!\mathrm{d}z\int_0^z\dfrac{y}{y^2+z^2}\,\mathrm{d}y
        =\dfrac{\ln2}{2}.
    \]
\end{solution}

\begin{example}
    计算
    \[
        \iiint\limits_V y\cos(x+z)\,\mathrm{d}x\,\mathrm{d}y\,\mathrm{d}z,
    \]
    其中 $V$ 是由 $y=\sqrt{x},x+z=\dfrac{\pi}{2},y=0,z=0$ 所围成的区域.
\end{example}
\exampleblank

\begin{solution}
    区域可写为 $0\leq x\leq\pi/2$,$0\leq y\leq\sqrt{x}$,$0\leq z\leq\pi/2-x$.于是
    \[
        I=\dfrac{1}{2}\int_0^{\frac{\pi}{2}}x(1-\sin x)\,\mathrm{d}x
        =\dfrac{\pi^2}{16}-\dfrac{1}{2}.
    \]
\end{solution}

\begin{example}
    计算
    \[
        \iiint\limits_V(xy+yz+zx)\,\mathrm{d}x\,\mathrm{d}y\,\mathrm{d}z,
    \]
    其中 $V:x\geq0,y\geq0,x^2+y^2\leq1,0\leq z\leq1$.
\end{example}
\exampleblank

\begin{solution}
    先对 $z$ 积分,并在第一象限单位圆上使用极坐标:
    \[
        I=\iint\limits_{\substack{x^2+y^2\leq1\\x,y\geq0}}
        \left[xy+\dfrac{x+y}{2}\right]\mathrm{d}x\,\mathrm{d}y
        =\dfrac{1}{8}+\dfrac{1}{3}=\dfrac{11}{24}.
    \]
\end{solution}



\subsection{变量替换}

\begin{theorem}[\markstar]
    设 $\widetilde\Omega,\Omega$ 是 $(x,y,z)$ 空间中的有界闭区域, 变换
    \[
        T:\begin{cases}
            x=x(u,v,w), \\
            y=y(u,v,w), \\
            z=z(u,v,w)
        \end{cases}
        :\widetilde\Omega\to\Omega
    \]
    是同胚变换, 且 $T\in C^{(2)}(\widetilde\Omega)$,
    \[
        J=\dfrac{\partial(x,y,z)}{\partial(u,v,w)}.
    \]
    若 $f(x,y,z)$ 在 $\Omega$ 上可积, 则
    \[
        \iiint\limits_\Omega f(x,y,z)\,\mathrm{d}x\,\mathrm{d}y\,\mathrm{d}z
        =\iiint\limits_{\widetilde\Omega}
        f\left(x(u,v,w),y(u,v,w),z(u,v,w)\right)|J|\,\mathrm{d}u\,\mathrm{d}v\,\mathrm{d}w.
    \]
\end{theorem}

\begin{remark}
    常用特例包括:
    \begin{enumerate}
        \item 正交变换
              \[
                  \begin{cases}
                      x=a_{11}u+a_{12}v+a_{13}w, \\
                      y=a_{21}u+a_{22}v+a_{23}w, \\
                      z=a_{31}u+a_{32}v+a_{33}w,
                  \end{cases}
              \]
              其中 $A=(a_{ij})$ 为正交矩阵, 此时 $|J|=1$;
        \item 柱坐标变换
              \[
                  x=r\cos\theta,\qquad y=r\sin\theta,\qquad z=z,
              \]
              其中 $0\leq r<+\infty,0\leq\theta<2\pi,-\infty<z<+\infty$, 此时 $J=r$;
        \item 球坐标变换
              \[
                  x=r\sin\varphi\cos\theta,\qquad
                  y=r\sin\varphi\sin\theta,\qquad
                  z=r\cos\varphi,
              \]
              其中 $0\leq r<+\infty,0\leq\varphi\leq\pi,0\leq\theta<2\pi$, 此时 $J=r^2\sin\varphi$.
    \end{enumerate}
\end{remark}

\begin{example}
    \begin{enumerate}
        \item 证明
              \[
                  \iint\limits_S f(ax+by)\,\mathrm{d}x\,\mathrm{d}y
                  =2\int_{-1}^1\sqrt{1-u^2}\,f\left(u\sqrt{a^2+b^2}\right)\mathrm{d}u,
              \]
              其中 $S:x^2+y^2\leq1,a^2+b^2\ne0$;
        \item 证明
              \[
                  \iiint\limits_Vf(ax+by+cz)\,\mathrm{d}x\,\mathrm{d}y\,\mathrm{d}z
                  =\pi\int_{-1}^1(1-u^2)f\left(u\sqrt{a^2+b^2+c^2}\right)\mathrm{d}u,
              \]
              其中 $V:x^2+y^2+z^2\leq1,a^2+b^2+c^2\ne0$.
    \end{enumerate}
\end{example}
\exampleblank

\begin{solution}
    分别作正交旋转,使线性函数成为第一坐标的倍数.单位圆在该坐标下仍为单位圆,其固定第一坐标 $u$ 的截线长度为 $2\sqrt{1-u^2}$,故
    \[
        \iint_Sf(ax+by)\,\mathrm{d}x\,\mathrm{d}y
        =2\int_{-1}^1\sqrt{1-u^2}\,f\left(u\sqrt{a^2+b^2}\right)\mathrm{d}u.
    \]
    同理,单位球固定第一坐标 $u$ 的截面面积为 $\pi(1-u^2)$,于是得到第二个等式.
\end{solution}

\begin{example}
    计算积分
    \[
        \iiint\limits_\Omega x^2\sqrt{x^2+y^2}\,\mathrm{d}x\,\mathrm{d}y\,\mathrm{d}z,
    \]
    其中 $\Omega$ 是曲面 $z=\sqrt{x^2+y^2}$ 与 $z=x^2+y^2$ 围成的有界闭区域.
\end{example}
\exampleblank

\begin{solution}
    用柱坐标.区域满足 $0\leq r\leq1$,$r^2\leq z\leq r$.故
    \[
        I=\int_0^{2\pi}\cos^2\theta\,\mathrm{d}\theta
        \int_0^1r^4(r-r^2)\,\mathrm{d}r
        =\pi\left(\dfrac{1}{6}-\dfrac{1}{7}\right)=\dfrac{\pi}{42}.
    \]
\end{solution}

\begin{example}
    \begin{enumerate}
        \item 计算椭圆
              \[
                  (a_1x+b_1y+c_1)^2+(a_2x+b_2y+c_2)^2=1,
                  \qquad a_1b_2-b_1a_2\ne0
              \]
              的面积;
        \item 计算椭球
              \[
                  (a_1x+b_1y+c_1z)^2+(a_2x+b_2y+c_2z)^2+(a_3x+b_3y+c_3z)^2=1
              \]
              的体积;
        \item 计算由平面
              \[
                  a_1x+b_1y+c_1z=\pm1,\quad
                  a_2x+b_2y+c_2z=\pm1,\quad
                  a_3x+b_3y+c_3z=\pm1
              \]
              所围立体的体积.
    \end{enumerate}
\end{example}
\exampleblank

\begin{solution}
    \begin{enumerate}
        \item 令 $u_i=a_ix+b_iy+c_i$,则单位圆面积除以线性变换行列式的绝对值,故
              \[
                  S=\dfrac{\pi}{|a_1b_2-b_1a_2|}.
              \]
        \item 记三阶系数矩阵为 $A$,则
              \[
                  V=\dfrac{4\pi}{3|\det A|}.
              \]
        \item 同样令三个平面左端为新坐标,像区域为立方体 $[-1,1]^3$,故
              \[
                  V=\dfrac{8}{|\det A|}.
              \]
    \end{enumerate}
\end{solution}

\begin{example}
    计算积分
    \[
        \iiint\limits_\Omega z^2\,\mathrm{d}x\,\mathrm{d}y\,\mathrm{d}z,
    \]
    其中 $\Omega$ 是 $x^2+y^2+z^2\leq a^2$ 与 $x^2+y^2+(z-a)^2\leq a^2$ $(a>0)$ 的公共部分.
\end{example}
\exampleblank

\begin{solution}
    用柱坐标.两个球的公共部分满足 $0\leq z\leq a$;在 $0\leq z\leq a/2$ 时截面半径平方为 $a^2-(z-a)^2$,在 $a/2\leq z\leq a$ 时为 $a^2-z^2$.因此
    \[
        I=\pi\int_0^{\frac{a}{2}}z^2\left[a^2-(z-a)^2\right]\,\mathrm{d}z
        +\pi\int_{\frac{a}{2}}^{a}z^2(a^2-z^2)\,\mathrm{d}z
        =\dfrac{59\pi a^5}{480}.
    \]
\end{solution}

\begin{example}
    计算积分
    \[
        \iiint\limits_{x^2+y^2+z^2\leq1}\cos(ax+by+cz)\,\mathrm{d}x\,\mathrm{d}y\,\mathrm{d}z.
    \]
\end{example}
\exampleblank

\begin{solution}
    令 $\lambda=\sqrt{a^2+b^2+c^2}$,作正交旋转使 $ax+by+cz=\lambda z$.球坐标积分给出
    \[
        I=4\pi\int_0^1r^2\dfrac{\sin(\lambda r)}{\lambda r}\,\mathrm{d}r
        =4\pi\dfrac{\sin\lambda-\lambda\cos\lambda}{\lambda^3}.
    \]
    当 $\lambda=0$ 时取连续延拓,值为 $4\pi/3$.
\end{solution}

\begin{example}
    计算积分:
    \begin{enumerate}
        \item $\displaystyle\iiint\limits_V(x^2+y^2)\,\mathrm{d}V$, 其中 $V$ 是由 $x^2+y^2+(z-2)^2\geq4$, $x^2+y^2+(z-1)^2\leq9$, $z\geq0$ 所围成的空心立体;
        \item
              \[
                  \iiint\limits_V\dfrac{xyz}{x^2+y^2}\,\mathrm{d}x\,\mathrm{d}y\,\mathrm{d}z,
              \]
              其中 $V$ 是曲面 $(x^2+y^2+z^2)^2=a^2xy$ $(a>0)$ 位于第一象限所围立体;
        \item
              \[
                  \iiint\limits_V(x+y+z)^2\,\mathrm{d}x\,\mathrm{d}y\,\mathrm{d}z,
              \]
              其中 $V$ 为 $x^2+y^2\leq2az$, $x^2+y^2+z^2\leq3$ 所围立体.
    \end{enumerate}
\end{example}
\exampleblank

\begin{solution}
    \begin{enumerate}
        \item 大球在 $z\geq0$ 内的部分减去内切小球,按水平截面积分可得
              \[
                  \dfrac{\pi}{2}\int_0^4\left(\left[9-(z-1)^2\right]^2-\left[4-(z-2)^2\right]^2\right)\mathrm{d}z
                  =\dfrac{256\pi}{3}.
              \]
        \item 球坐标中边界为 $\rho^2=a^2\sin^2\varphi\cos\theta\sin\theta$.在第一卦限积分得到
              \[
                  I=\dfrac{a^4}{4}\int_0^{\frac{\pi}{2}}\sin^5\varphi\cos\varphi\,\mathrm{d}\varphi
                  \int_0^{\frac{\pi}{2}}(\sin\theta\cos\theta)^3\,\mathrm{d}\theta
                  =\dfrac{a^4}{288}.
              \]
        \item 令 $\varphi_0\in(0,\pi/2)$ 满足
              \[
                  \dfrac{2a\cos\varphi_0}{\sin^2\varphi_0}=\sqrt3,
                  \qquad u_0=\cot\varphi_0.
              \]
              对方位角积分后,原积分为
              \[
                  \dfrac{2\pi}{5}\left[3^{\frac{5}{2}}(1-\cos\varphi_0)
                      +32a^5\left(\dfrac{u_0^6}{6}+\dfrac{u_0^8}{8}\right)\right].
              \]
    \end{enumerate}
\end{solution}

\begin{example}[\markstar]
    \begin{enumerate}
        \item 设 $a,b,c>0$, 证明
              \[
                  \iiint\limits_Vx^{a-1}y^{b-1}z^{c-1}\,\mathrm{d}x\,\mathrm{d}y\,\mathrm{d}z
                  =\dfrac{1}{a+b+c}\dfrac{\Gamma(a)\Gamma(b)\Gamma(c)}{\Gamma(a+b+c)},
              \]
              其中 $V$ 为四面体 $x\geq0,y\geq0,z\geq0,x+y+z\leq1$;
        \item 计算
              \[
                  \iiint\limits_Vxy^9z^8w^4\,\mathrm{d}x\,\mathrm{d}y\,\mathrm{d}z,
                  \qquad w=1-x-y-z,
              \]
              其中 $V:x,y,z\geq0,x+y+z\leq1$.
    \end{enumerate}
\end{example}
\exampleblank

\begin{solution}
    \begin{enumerate}
        \item 逐次使用 Beta 积分即得
              \[
                  \iiint_Vx^{a-1}y^{b-1}z^{c-1}\,\mathrm{d}V
                  =\dfrac{\Gamma(a)\Gamma(b)\Gamma(c)}{\Gamma(a+b+c+1)}
                  =\dfrac{1}{a+b+c}\dfrac{\Gamma(a)\Gamma(b)\Gamma(c)}{\Gamma(a+b+c)}.
              \]
        \item 把 $w=1-x-y-z$ 看作第四个单纯形坐标,Dirichlet 积分给出
              \[
                  \iiint_Vxy^9z^8w^4\,\mathrm{d}V
                  =\dfrac{\Gamma(2)\Gamma(10)\Gamma(9)\Gamma(5)}{\Gamma(26)}
                  =\dfrac{9!\,8!\,4!}{25!}.
              \]
    \end{enumerate}
\end{solution}



\subsection{理论及应用}

\begin{example}
    设 $f(u)$ 在 $\lbrack0,+\infty)$ 上可微, 试求极限
    \[
        \lim\limits_{t\to0^+}\dfrac{1}{\pi t^4}
        \iiint\limits_{x^2+y^2+z^2\leq t^2}
        f\left(\sqrt{x^2+y^2+z^2}\right)\,\mathrm{d}x\,\mathrm{d}y\,\mathrm{d}z.
    \]
\end{example}
\exampleblank

\begin{solution}
    球坐标给出
    \[
        \dfrac{1}{\pi t^4}\iiint_{r\leq t}f(r)\,\mathrm{d}V
        =\dfrac{4}{t^4}\int_0^t f(r)r^2\,\mathrm{d}r.
    \]
    若 $f(0)=0$,则 $f(r)=f'(0)r+o(r)$,极限为 $f'(0)$.若 $f(0)\ne0$,上式主项为 $4f(0)/(3t)$,不存在有限极限.因此原题若要求有限答案,需补充 $f(0)=0$.
\end{solution}

\begin{example}[\markstar]
    证明
    \[
        \lim\limits_{n\to\infty}\dfrac{1}{n^4}
        \iiint\limits_{r\leq n}[r]\,\mathrm{d}x\,\mathrm{d}y\,\mathrm{d}z=\pi,
        \qquad r=\sqrt{x^2+y^2+z^2}.
    \]
\end{example}
\exampleblank

\begin{solution}
    由球坐标及 $[r]=r+O(1)$,
    \[
        \dfrac{1}{n^4}\iiint_{r\leq n}[r]\,\mathrm{d}V
        =\dfrac{4\pi}{n^4}\int_0^n[r]r^2\,\mathrm{d}r
        =\dfrac{4\pi}{n^4}\left(\dfrac{n^4}{4}+O(n^3)\right)\to\pi.
    \]
\end{solution}

\begin{example}
    设 $f(x,y,z)$ 在
    \[
        V=\{(x,y,z)\mid0\leq x,y,z\leq1\}
    \]
    上有六阶连续偏导数, $f$ 在边界上恒为零, 且
    \[
        \left|\dfrac{\partial^6f(x,y,z)}{\partial x^2\partial y^2\partial z^2}\right|\leq M,
    \]
    证明
    \[
        \left|\iiint\limits_Vf(x,y,z)\,\mathrm{d}x\,\mathrm{d}y\,\mathrm{d}z\right|
        \leq\dfrac{1}{8}\cdot\dfrac{M}{6^3}.
    \]
\end{example}
\exampleblank

\begin{solution}
    反复使用一维恒等式
    \[
        \int_0^1g(t)\,\mathrm{d}t=-\dfrac{1}{2}\int_0^1t(1-t)g''(t)\,\mathrm{d}t
    \]
    (其中 $g(0)=g(1)=0$),得到
    \[
        \iiint_Vf=-\dfrac{1}{8}\iiint_Vx(1-x)y(1-y)z(1-z)f_{xxyyzz}\,\mathrm{d}V.
    \]
    取绝对值并用 $\int_0^1t(1-t)\,\mathrm{d}t=1/6$,即得
    \[
        \left|\iiint_Vf\right|\leq\dfrac{1}{8}\cdot\dfrac{M}{6^3}.
    \]
\end{solution}

\begin{example}
    求由下列曲面围成的立体的体积 $V$:
    \begin{enumerate}
        \item $\left(\dfrac{x^2}{a^2}+\dfrac{y^2}{b^2}+\dfrac{z^2}{c^2}\right)^2=hx$;
        \item $(x^2+y^2)^2+z^4=y$;
        \item $x^2+z^2=a^2$, $x^2+z^2=b^2$, $x^2-y^2-z^2=0$, $x>0$, $a<b$.
    \end{enumerate}
\end{example}
\exampleblank

\begin{solution}
    \begin{enumerate}
        \item 令 $x=aX,y=bY,z=cZ$,再以 $X$ 轴为球坐标极轴.边界满足 $R^3=ah\cos\theta$,故
              \[
                  V=abc\cdot\dfrac{1}{3}ah\int_{\cos\theta\geq0}\cos\theta\,\mathrm{d}\Omega
                  =\dfrac{\pi a^2bch}{3}.
              \]
        \item 柱坐标中先对 $z$ 积分,并令 $t=r^3/\sin\theta$,得到
              \[
                  V=\dfrac{4}{3} B\left(\dfrac{3}{4},\dfrac{5}{4}\right)=\dfrac{\pi\sqrt2}{3}.
              \]
        \item 以 $y$ 轴为轴作柱坐标,$a\leq r\leq b$,$-r\leq y\leq r$,而 $x>0$ 使方位角长度为 $\pi$.所以
              \[
                  V=\pi\int_a^b2r^2\,\mathrm{d}r=\dfrac{2\pi}{3}(b^3-a^3).
              \]
    \end{enumerate}
\end{solution}

\begin{example}[\markstar]
    设
    \[
        \sum\limits_{i,j=1}^3a_{ij}x_ix_j
    \]
    表示变量 $(x_1,x_2,x_3)$ 的二次型, 其中 $A=(a_{ij})$ 正定. 证明椭球面
    \[
        S:\quad\sum\limits_{i,j=1}^3a_{ij}x_ix_j=1
    \]
    所包围的体积等于
    \[
        \dfrac{4}{3}\pi(\det A)^{-\frac{1}{2}}.
    \]
\end{example}
\exampleblank

\begin{solution}
    取正定平方根 $A^{1/2}$ 并令 $\bm y=A^{1/2}\bm x$.则 $S$ 的内部变成单位球,且
    \[
        \mathrm{d}\bm x=(\det A)^{-\frac{1}{2}}\mathrm{d}\bm y.
    \]
    因此
    \[
        V=(\det A)^{-\frac{1}{2}}\cdot\dfrac{4\pi}{3}.
    \]
\end{solution}



\subsection{\texorpdfstring{$n$}{n} 重积分 *}

$n$ 重积分处理技巧基本上类似于二,三重积分, 由于 $n$ 维的一般性, 计算时多了一种递推方法.

\begin{theorem}[\markstar\ 化为累次积分]
    \begin{enumerate}
        \item 设
              \[
                  V=\{(x_1,x_2,\ldots,x_n)\mid a_i\leq x_i\leq b_i,\ i=1,2,\ldots,n\},
              \]
              $f$ 在 $V$ 上连续, 则
              \[
                  \int\cdots\int\limits_V f(x_1,\ldots,x_n)\,\mathrm{d}x_1\cdots\mathrm{d}x_n
                  =\int_{a_1}^{b_1}\mathrm{d}x_1\cdots\int_{a_n}^{b_n}f(x_1,\ldots,x_n)\,\mathrm{d}x_n.
              \]

        \item 设
              \[
                  \begin{aligned}
                      V=\{(x_1,\ldots,x_n)\mid{} & a_1\leq x_1\leq b_1,\ a_2(x_1)\leq x_2\leq b_2(x_1),\ldots,    \\
                                                 & a_n(x_1,\ldots,x_{n-1})\leq x_n\leq b_n(x_1,\ldots,x_{n-1})\},
                  \end{aligned}
              \]
              $f$ 在 $V$ 上连续, 则相应的 $n$ 重积分可按上述边界依次化为累次积分.
    \end{enumerate}
\end{theorem}

换元结论类似. 值得一提的是推广的球坐标变换:
\[
    \begin{cases}
        x_1=r\cos\varphi_1,                                             \\
        x_2=r\sin\varphi_1\cos\varphi_2,                                \\
        x_3=r\sin\varphi_1\sin\varphi_2\cos\varphi_3,                   \\
        \cdots,                                                         \\
        x_{n-1}=r\sin\varphi_1\cdots\sin\varphi_{n-2}\cos\varphi_{n-1}, \\
        x_n=r\sin\varphi_1\cdots\sin\varphi_{n-2}\sin\varphi_{n-1}.
    \end{cases}
\]
其中 $0\leq r<+\infty$, $0\leq\varphi_1,\ldots,\varphi_{n-2}\leq\pi$, $0\leq\varphi_{n-1}<2\pi$, Jacobi 行列式为
\[
    J=r^{n-1}\sin^{n-2}\varphi_1\sin^{n-3}\varphi_2\cdots\sin^2\varphi_{n-3}\sin\varphi_{n-2}.
\]

\begin{example}
    设 $a_1,\ldots,a_n>0$, 又
    \[
        S_n(a_1,\ldots,a_n)=\left\{(x_1,\ldots,x_n)\ \middle|\ \dfrac{|x_1|}{a_1}+\cdots+\dfrac{|x_n|}{a_n}\leq1\right\},
    \]
    计算 $S_n(a_1,\ldots,a_n)$ 的体积.
\end{example}
\exampleblank

\begin{solution}
    先看第一卦限,其中令 $u_i=x_i/a_i$ 后得到标准单纯形,体积为 $(a_1\cdots a_n)/n!$.共有 $2^n$ 个全等卦限,故
    \[
        \operatorname{Vol}S_n(a_1,\ldots,a_n)=\dfrac{2^n}{n!}\prod_{i=1}^na_i.
    \]
\end{solution}

\begin{example}
    求
    \[
        \iiiint\limits_V\sqrt{\dfrac{1-x^2-y^2-z^2-u^2}{1+x^2+y^2+z^2+u^2}}\,
        \mathrm{d}x\,\mathrm{d}y\,\mathrm{d}z\,\mathrm{d}u,
    \]
    其中 $V:x,y,z,u\geq0$, $x^2+y^2+z^2+u^2\leq1$.
\end{example}
\exampleblank

\begin{solution}
    四维单位球面面积为 $2\pi^2$,第一卦限占 $1/16$.故原积分为
    \[
        \dfrac{\pi^2}{8}\int_0^1r^3\sqrt{\dfrac{1-r^2}{1+r^2}}\,\mathrm{d}r.
    \]
    令 $t=r^2$,再令 $u=\sqrt{(1-t)/(1+t)}$,可得
    \[
        \int_0^1r^3\sqrt{\dfrac{1-r^2}{1+r^2}}\,\mathrm{d}r
        =\dfrac{1}{2}\left(1-\dfrac{\pi}{4}\right).
    \]
    所以结果为
    \[
        \dfrac{\pi^2(4-\pi)}{64}.
    \]
\end{solution}

\begin{example}
    计算
    \[
        \iiiint\limits_De^{-\langle A\bm{x},\bm{x}\rangle}\,
        \mathrm{d}x_1\,\mathrm{d}x_2\,\mathrm{d}x_3\,\mathrm{d}x_4,
    \]
    其中
    \[
        \langle A\bm{x},\bm{x}\rangle=\sum\limits_{i,j=1}^4a_{ij}x_ix_j
    \]
    是正定二次型, $D$ 是区域 $\langle A\bm{x},\bm{x}\rangle\leq1$.
\end{example}
\exampleblank

\begin{solution}
    令 $\bm y=A^{1/2}\bm x$,则 $D$ 变成四维单位球,Jacobian 为 $(\det A)^{-1/2}$.四维球坐标给出
    \[
        I=\dfrac{2\pi^2}{\sqrt{\det A}}\int_0^1e^{-r^2}r^3\,\mathrm{d}r
        =\dfrac{\pi^2}{\sqrt{\det A}}\left(1-\dfrac{2}{e}\right).
    \]
\end{solution}

\begin{example}
    试求 $\mathbb{R}^n$ 中下列立体 $\Omega$ 的体积 $V$:
    \begin{enumerate}
        \item
              \[
                  \Omega_n(a)=\{(x_1,\ldots,x_n)\mid x_i\geq0,\ x_1+\cdots+x_n\leq a\};
              \]
        \item
              \[
                  \Omega=\left\{(x_1,\ldots,x_n)\ \middle|\ x_1^{a_1}+\cdots+x_n^{a_n}\leq1,\ x_i\geq0\right\},
              \]
              其中 $a_i>0$;
        \item
              \[
                  \Omega=\{(x_1,\ldots,x_n)\mid x_1^2+x_2^2+\cdots+x_n^2\leq a^2\},\qquad a>0.
              \]
    \end{enumerate}
\end{example}
\exampleblank

\begin{solution}
    \begin{enumerate}
        \item 标准单纯形的缩放直接给出
              \[
                  V=\dfrac{a^n}{n!}.
              \]
        \item 令 $u_i=x_i^{a_i}$,由 Dirichlet 积分得到
              \[
                  V=\dfrac{\displaystyle\prod_{i=1}^n\Gamma\left(1+\dfrac{1}{a_i}\right)}
                  {\displaystyle\Gamma\left(1+\sum_{i=1}^n\dfrac{1}{a_i}\right)}.
              \]
        \item 设单位 $n$ 维球体积为 $V_n$.由高斯积分 $\pi^{n/2}=V_n\Gamma(n/2+1)$,故
              \[
                  V=\dfrac{\pi^{n/2}a^n}{\Gamma\left(\dfrac{n}{2}+1\right)}.
              \]
    \end{enumerate}
\end{solution}

\begin{example}
    设 $f(x)$ 在 $\lbrack a,b\rbrack$ 上连续, 证明: 对任意 $x\in(a,b)$,
    \[
        \int_a^x\mathrm{d}x_1\int_a^{x_1}\mathrm{d}x_2\cdots\int_a^{x_n}f(x_{n+1})\,\mathrm{d}x_{n+1}
        =\dfrac{1}{n!}\int_a^x(x-y)^nf(y)\,\mathrm{d}y.
    \]
\end{example}
\exampleblank

\begin{solution}
    对 $n$ 归纳.$n=0$ 时即微积分基本定理.假设结论对 $n-1$ 成立,则最外层再积分一次并交换三角形区域次序:
    \[
        \begin{aligned}
            \int_a^x\dfrac{1}{(n-1)!}\int_a^{x_1}(x_1-y)^{n-1}f(y)\,\mathrm{d}y\,\mathrm{d}x_1
             & =\dfrac{1}{(n-1)!}\int_a^xf(y)\int_y^x(x_1-y)^{n-1}\,\mathrm{d}x_1\,\mathrm{d}y \\
             & =\dfrac{1}{n!}\int_a^x(x-y)^nf(y)\,\mathrm{d}y.
        \end{aligned}
    \]
\end{solution}

\begin{example}
    设 $f\in C\lbrack0,1\rbrack$, 且
    \[
        \int_0^1f(x)\,\mathrm{d}x=A,
    \]
    求积分
    \[
        \int_0^1\cdots\int_0^1
        \dfrac{x_1+x_2+\cdots+x_k}{x_1+x_2+\cdots+x_n}
        f(x_1)f(x_2)\cdots f(x_n)
        \,\mathrm{d}x_1\cdots\mathrm{d}x_n,
        \qquad k\leq n.
    \]
\end{example}
\exampleblank

\begin{solution}
    记所求积分中分子为 $x_j$ 时的值为 $I_j$.由变量的置换对称性,$I_1=\cdots=I_n$,且
    \[
        \sum_{j=1}^nI_j=\int_{[0,1]^n}\prod_{j=1}^nf(x_j)\,\mathrm{d}x_1\cdots\mathrm{d}x_n=A^n.
    \]
    所以原积分为
    \[
        I_1+\cdots+I_k=\dfrac{k}{n}A^n.
    \]
\end{solution}


\newpage

\section{反常重积分 *}

\begin{definition}[无界区域]
    设 $D$ 是 $\mathbb{R}^2$ 中的无界区域, 其边界由有限条光滑或逐段光滑曲线组成. 函数 $f$ 定义在 $D$ 上, 且在 $D$ 内的任何可求面积的有界子区域上可积. 设 $D_r$ 是 $D$ 内的任一可求面积的有界子区域, 包含 $D\cap B_r$, 其中 $B_r$ 是 $\mathbb{R}^2$ 中以 $r$ 为半径的闭圆盘. 若极限
    \[
        I=\lim\limits_{r\to+\infty}\iint\limits_{D_r}f(x,y)\,\mathrm{d}x\,\mathrm{d}y
    \]
    存在且有限, 并与 $D_r$ 的取法无关, 则称 $f$ 在 $D$ 上积分收敛, 否则称积分发散.
\end{definition}

\begin{remark}
    因此, 证明发散的话, 可以找两种收敛于不同值的分割方法, 或一种极限不存在的分割方式.
\end{remark}

\begin{theorem}[\markstar]
    设 $f(x,y)$ 是无界区域 $D\subseteq\mathbb{R}^2$ 上的非负函数, 且在 $D$ 内的任何可求面积的有界子区域上可积, $D_n$ 是包含 $D\cap B_n$ 的一列可求面积的有界子区域, 则
    \[
        \iint\limits_D f(x,y)\,\mathrm{d}x\,\mathrm{d}y\text{ 收敛}
        \iff
        \left\{\iint\limits_{D_n}f(x,y)\,\mathrm{d}x\,\mathrm{d}y\right\}\text{ 有界}
    \]
    且等价于
    \[
        \lim\limits_{n\to+\infty}\iint\limits_{D_n}f(x,y)\,\mathrm{d}x\,\mathrm{d}y
    \]
    存在.
\end{theorem}

\begin{remark}
    即非负函数时, 只需找到一种特殊的区间套, 使得 $f(x,y)$ 在其上积分有界即可.
\end{remark}

\begin{definition}[无界函数]
    设 $D$ 为 $\mathbb{R}^2$ 上可求面积的有界区域, 点 $P_0\in D$, 函数 $f$ 定义在 $D\setminus\{P_0\}$ 上且对任何点 $P_0$ 的可求面积的邻域 $\Delta$, $f$ 在 $D\setminus\Delta$ 上有界可积. 如果
    \[
        I=\lim\limits_{d(\Delta)\to0}\iint\limits_{D\setminus\Delta}f(x,y)\,\mathrm{d}x\,\mathrm{d}y
    \]
    存在且有界, 并与 $\Delta$ 的取法无关, 其中 $d(\Delta)$ 为 $\Delta$ 的直径, 则称 $f$ 在 $D$ 上积分收敛, 否则称积分发散.
\end{definition}

\begin{remark}
    显然, 无界函数也有相应于前述定理的非负函数充要条件版本. 类似于一元版本, 对于非负函数有相应比较判别法. 特别地, 当取 $g(x,y)=\dfrac{C}{r^p}$, $r=\sqrt{x^2+y^2}$ 时, 可得到下面的 Cauchy 判别法.
\end{remark}

\begin{theorem}[\markstar\ Cauchy 判别法]
    设 $D$ 为无界区域, 其边界由有限条光滑或逐段连续曲线组成, $f$ 在 $D$ 内的任一可求面积的有界子区域上可积, 则:
    \begin{enumerate}
        \item 若对充分大的 $r$, 有 $|f(x,y)|\leq\dfrac{C}{r^p}$, $p>2$, 则反常积分 $\displaystyle\iint\limits_Df(x,y)\,\mathrm{d}x\,\mathrm{d}y$ 收敛;
        \item 若对充分大的 $r$, 有 $|f(x,y)|>\dfrac{C}{r^p}$, $p\leq2$, 则反常积分 $\displaystyle\iint\limits_Df(x,y)\,\mathrm{d}x\,\mathrm{d}y$ 发散.
    \end{enumerate}
\end{theorem}

总的来说, 反常重积分与反常积分大部分理论完全相似, 但是却有如下惊人的差别.

\begin{theorem}[\markstar]
    $f$ 反常可积当且仅当 $|f|$ 反常可积.
\end{theorem}

\begin{example}
    讨论反常重积分
    \[
        I=\iint\limits_D\dfrac{\mathrm{d}x\,\mathrm{d}y}{(x+y)^p}
    \]
    的敛散性, 其中 $D=\{(x,y)\mid0\leq x\leq1,\ x+y\geq1\}$. 收敛时, 求积分的值.
\end{example}
\exampleblank

\begin{solution}
    令 $u=x+y$.对每个 $x\in[0,1]$,$u$ 从 $1$ 到 $+\infty$,故
    \[
        I=\int_0^1\!\mathrm{d}x\int_1^{+\infty}u^{-p}\,\mathrm{d}u.
    \]
    因此积分当且仅当 $p>1$ 时收敛,并且此时
    \[
        I=\dfrac{1}{p-1}.
    \]
\end{solution}

\begin{example}
    设 $D:1\leq x,y<+\infty$,
    \[
        f(x,y)=\dfrac{x^2-y^2}{(x^2+y^2)^2},
    \]
    令
    \[
        I=\iint\limits_Df(x,y)\,\mathrm{d}x\,\mathrm{d}y,
    \]
    \[
        I_1=\int_1^{+\infty}\mathrm{d}x\int_1^{+\infty}f(x,y)\,\mathrm{d}y,
        \qquad
        I_2=\int_1^{+\infty}\mathrm{d}y\int_1^{+\infty}f(x,y)\,\mathrm{d}x.
    \]
    则 $I$ 发散, $I_1$ 与 $I_2$ 收敛.
\end{example}
\exampleblank

\begin{solution}
    利用
    \[
        \dfrac{x^2-y^2}{(x^2+y^2)^2}
        =\dfrac{\partial}{\partial y}\dfrac{y}{x^2+y^2}
        =-\dfrac{\partial}{\partial x}\dfrac{x}{x^2+y^2},
    \]
    可得
    \[
        I_1=-\int_1^{+\infty}\dfrac{\mathrm{d}x}{1+x^2}=-\dfrac{\pi}{4},
        \qquad
        I_2=\int_1^{+\infty}\dfrac{\mathrm{d}y}{1+y^2}=\dfrac{\pi}{4}.
    \]
    若二重反常积分存在,Fubini 型结论会迫使两者相等,矛盾;故 $I$ 发散.
\end{solution}

\begin{example}
    证明积分
    \[
        I=\iint\limits_D\sin(x^2+y^2)\,\mathrm{d}x\,\mathrm{d}y
    \]
    不存在, 其中 $D:0\leq x,y<+\infty$.
\end{example}
\exampleblank

\begin{solution}
    取四分之一圆盘 $D_R$ 截断并使用极坐标,则
    \[
        \iint_{D_R}\sin(x^2+y^2)\,\mathrm{d}x\,\mathrm{d}y
        =\dfrac{\pi}{4}\int_0^{R^2}\sin u\,\mathrm{d}u
        =\dfrac{\pi}{4}(1-\cos R^2).
    \]
    当 $R\to+\infty$ 时该值没有极限,因此题给反常积分不存在.
\end{solution}

\begin{example}
    计算下列反常积分:
    \begin{enumerate}
        \item $\displaystyle I=\iint\limits_D\dfrac{\mathrm{d}x\,\mathrm{d}y}{(1+x+y)^3}$, $D:x\geq0,y\geq0$;
        \item $\displaystyle I=\iint\limits_D\dfrac{\mathrm{d}x\,\mathrm{d}y}{(x+y)^p}$, $D:0\leq x\leq1,x+y\geq1$.
    \end{enumerate}
\end{example}
\exampleblank

\begin{solution}
    \begin{enumerate}
        \item 令 $u=x+y$,同一直线截段长度为 $u$,故
              \[
                  I=\int_0^{+\infty}\dfrac{u}{(1+u)^3}\,\mathrm{d}u=\dfrac{1}{2}.
              \]
        \item 与上一题相同,$p>1$ 时收敛且 $I=1/(p-1)$;$p\leq1$ 时发散.
    \end{enumerate}
\end{solution}

\begin{example}
    判别下列积分的敛散性:
    \begin{enumerate}
        \item
              \[
                  I=\iint\limits_D\dfrac{\varphi(x,y)}{(x^2+y^2)^p}\,\mathrm{d}x\,\mathrm{d}y,
              \]
              其中 $D:x^2+y^2\geq1,x\geq0,y\geq0$, $\varphi\in C(D)$ 且 $m\leq|\varphi(x,y)|\leq M$;
        \item
              \[
                  I=\iint\limits_D\dfrac{\mathrm{d}x\,\mathrm{d}y}{x^p+x^q},
                  \qquad p,q>0,
              \]
              其中 $D:x\geq0,y\geq0,x+y\geq1$;
        \item
              \[
                  I=\iint\limits_D\dfrac{\sin x\sin y}{(x+y)^p}\,\mathrm{d}x\,\mathrm{d}y,
              \]
              其中 $D:x+y\geq1$.
    \end{enumerate}
\end{example}
\exampleblank

\begin{solution}
    \begin{enumerate}
        \item 由 $m\leq|\varphi|\leq M$,绝对收敛性等价于
              \[
                  \int_1^{+\infty}r^{1-2p}\,\mathrm{d}r<+\infty,
              \]
              即 $p>1$.若只问非绝对收敛,则现有条件不能控制 $\varphi$ 的符号振荡,无法作统一判断.
        \item 按现有题面,被积函数与 $y$ 无关,而对每个允许的 $x$,$y$ 都延伸到 $+\infty$,故对一切 $p,q>0$ 均发散.若原题预期其他判据,分母很可能漏写了含 $y$ 的项.
        \item 若按通常意图补充 $x,y\geq0$,令 $u=x+y$,则内层积分为
              \[
                  \int_0^u\sin x\sin(u-x)\,\mathrm{d}x=\dfrac{1}{2}(\sin u-u\cos u).
              \]
              因而积分在 $p>1$ 时收敛,在 $p\leq1$ 时发散;其中 $p>2$ 时绝对收敛,$1<p\leq2$ 时条件收敛.原题若不补充第一象限条件,区域还穿过 $x+y=0$,题意不完整.
    \end{enumerate}
\end{solution}

\begin{example}
    设 $D:0\leq x\leq1,0\leq y\leq1,1\leq p<+\infty$, 令
    \[
        f(x,y)=
        \begin{cases}
            0,               & xy=1,   \\
            \dfrac{1}{xy-1}, & xy\ne1,
        \end{cases}
    \]
    则
    \[
        \iint\limits_D|f(x,y)|^p\,\mathrm{d}x\,\mathrm{d}y\text{ 收敛}
        \iff1\leq p\leq2.
    \]
\end{example}
\exampleblank

\begin{solution}
    在奇点 $(1,1)$ 附近令 $u=1-x$,$v=1-y$,则
    \[
        1-xy=u+v-uv\asymp u+v.
    \]
    因此局部可积性等价于
    \[
        \int_0^\varepsilon\!\int_0^\varepsilon(u+v)^{-p}\,\mathrm{d}u\,\mathrm{d}v<+\infty,
    \]
    其充要条件为 $p<2$.所以在题设 $p\geq1$ 下,正确结论是
    \[
        \iint_D|f|^p<+\infty\iff1\leq p<2.
    \]
    题面把 $p=2$ 误列为收敛端点;该端点实际为对数发散.
\end{solution}

\begin{example}
    讨论积分
    \[
        I=\iiint\limits_D\dfrac{\mathrm{d}x\,\mathrm{d}y\,\mathrm{d}z}{|x|^p+|y|^q+|z|^r},
        \qquad p,q,r>0,
    \]
    的敛散性, 其中 $D:|x|+|y|+|z|\geq1$.
\end{example}
\exampleblank

\begin{solution}
    记
    \[
        s=\dfrac{1}{p}+\dfrac{1}{q}+\dfrac{1}{r}.
    \]
    作各向异性伸缩 $x=t^{1/p}u$,$y=t^{1/q}v$,$z=t^{1/r}w$.集合
    \[
        \{|x|^p+|y|^q+|z|^r\leq t\}
    \]
    的体积等于常数乘以 $t^s$.按层积分,积分在无穷远处与
    \[
        \int_1^{+\infty}t^{s-2}\,\mathrm{d}t
    \]
    同敛散.区域已避开原点,故没有有限点奇性,于是
    \[
        I\text{ 收敛}\iff \dfrac{1}{p}+\dfrac{1}{q}+\dfrac{1}{r}<1.
    \]
\end{solution}
